%-------------------------------------------------------------------------------
%	ABSTRACT PAGE
%-------------------------------------------------------------------------------

\addchaptertocentry{Хураангуй}

\begin{center}
{\scshape\Large \univname\par}
{\scshape\large \facname\par}\vspace{0.5cm}
{\huge\textbf{{Хураангуй}} \par}
\bigskip
{\Large{\ttitle} \par}
\bigskip
{\normalsize \shortname \par}
\addressname
\end{center}

\textit{\textbf{Түлхүүр үгс: \keywordnames}}
\bigskip

Энэхүү дипломын ажлын хүрээнд хэрэглэгч зөвхөн сэдэв оруулахад скрипт бичих, дүрслэл үүсгэх, дуу хоолой синтез хийх, видео нэгтгэн хадгалах бүх үйл явцыг автоматаар гүйцэтгэдэг хиймэл оюунд суурилсан богино хэмжээний видео контент үүсгэлтийн цогц веб систем хөгжүүлэхэд чиглэсэн судалгаа, шинжилгээ болон системийн зохиомжийг боловсруулсан. \\

Системийн судалгааны хэсэгт InVideo AI, Pictory, Synthesia, Runway Gen-3 зэрэг олон улсын ижил төстэй системүүдтэй харьцуулалт хийж, манай системийн давуу тал болох Монгол хэлний дэмжлэгийг тодорхойлсон. Ашиглах AI загваруудын хувьд скрипт үүсгэлтэд Groq LLM (Llama 4), дуу хоолой синтезд ElevenLabs, CHIMEGE, дүрслэл үүсгэлтэд Google Gemini Flash Image 2.5 API-г судалж, техникийн үзүүлэлтийг нотлов. \\

Системийн шаардлагын хэсэгт гурван үүрэгтэй хэрэглэгч --- Хэрэглэгч, Контент Менежер, Супер Админ --- нарын функциональ болон функциональ бус шаардлагуудыг тодорхойлж, Use Case диаграм болон тодорхойлолтуудыг гаргав. Техникийн хэрэгжүүлэлтэд Node.js дээр суурилсан REST API backend, TypeScript, PostgreSQL өгөгдлийн сан, JWT баталгаажуулалт, FFmpeg видео боловсруулалт, Cloudinary/AWS S3 үүлэн хадгалалтыг ашиглах шийдвэр гарсан.
